\documentclass[11pt,a4paper]{article}
\usepackage[utf8]{inputenc}
\usepackage[T1]{fontenc}
\usepackage{lmodern}
\usepackage[french]{babel}
\usepackage{amsmath,amssymb,mathtools}
\usepackage{icomma}
\usepackage{xcolor}
\usepackage[most]{tcolorbox}
\usepackage{enumitem}
\usepackage[a4paper,margin=2.2cm,headheight=26pt,headsep=10pt]{geometry}
\usepackage{fancyhdr}
\usepackage[hidelinks]{hyperref}
\definecolor{primary}{RGB}{222,49,99}
\definecolor{primarydark}{RGB}{150,22,55}
\definecolor{excol}{RGB}{0,135,90}
\tcbset{boxbase/.style={enhanced,breakable,boxrule=0.9pt,arc=2.5pt,
  left=3mm,right=3mm,top=2.3mm,bottom=2.3mm,fonttitle=\bfseries,coltitle=white}}
\newtcolorbox[auto counter]{corrige}[1][]{boxbase,colback=excol!4,colframe=excol,
  title={\textsc{Corrigé}~\thetcbcounter\ifx\relax#1\relax\relax\else\textmd{ --- \itshape #1}\fi}}
\newcommand{\R}{\mathbb{R}}
\newcommand{\N}{\mathbb{N}}
\newcommand{\ds}{\displaystyle}
\pagestyle{fancy}\fancyhf{}
\fancyhead[L]{\small\textcolor{primarydark}{\textbf{Thème 1} — Puissances, racines et exposants}}
\fancyhead[R]{\small\textcolor{primarydark}{\textsc{Corrigé — Moyen}}}
\fancyfoot[C]{\small\thepage}
\renewcommand{\headrulewidth}{0.8pt}
\begin{document}

\begin{center}
{\LARGE\bfseries\color{primarydark} Niveau Moyen — Corrigé}\\[2pt]
{\color{primary}\rule{0.45\textwidth}{1.2pt}}
\end{center}
\vspace{1mm}

\begin{corrige}[combiner plusieurs règles]
\begin{align*}
\text{a) }&\frac{a^{5}a^{-3}}{a^{2}}=a^{5-3-2}=a^{0}=1; &
\text{b) }&(a^{2})^{3}a^{-4}=a^{6}a^{-4}=a^{2};\\
\text{c) }&\frac{(a^{3})^{2}}{a^{-1}}=\frac{a^{6}}{a^{-1}}=a^{6-(-1)}=a^{7}; &
\text{d) }&\Bigl(\frac{a^{-2}}{a^{3}}\Bigr)^{-1}=\bigl(a^{-5}\bigr)^{-1}=a^{5}.
\end{align*}
\end{corrige}

\begin{corrige}[racines et exposants fractionnaires]
\begin{align*}
\text{a) }&\frac{c^{3/2}}{c^{1/4}}=c^{3/2-1/4}=c^{5/4}; &
\text{b) }&c^{1/3}c^{1/6}=c^{2/6+1/6}=c^{1/2}=\sqrt{c};\\
\text{c) }&\Bigl((c^{2n^{2}})^{1/n}\Bigr)^{1/n}=c^{2n^{2}/n^{2}}=c^{2}; &
\text{d) }&(c^{1/2})^{6}c^{-2}=c^{3}c^{-2}=c^{1}=c.
\end{align*}
\end{corrige}

\begin{corrige}[bases opposées et quotients]
\[ \text{a) } a^{-m}b^{m}=\frac{b^{m}}{a^{m}}=\Bigl(\frac{b}{a}\Bigr)^{m};\quad
   \text{b) } \Bigl(\frac{b}{a}\Bigr)^{3};\quad
   \text{c) } 2^{-x}3^{x}=\Bigl(\frac{3}{2}\Bigr)^{x};\quad
   \text{d) } a^{4}b^{-4}=\Bigl(\frac{a}{b}\Bigr)^{4}. \]
\end{corrige}

\begin{corrige}[ramener à une base commune]
\begin{align*}
\text{a) }&4^{x}=2^{2x}=2^{6}\Rightarrow 2x=6\Rightarrow x=3; &
\text{b) }&2^{3x}=2^{x+4}\Rightarrow 3x=x+4\Rightarrow x=2;\\
\text{c) }&3^{2x}=3^{-3}\Rightarrow 2x=-3\Rightarrow x=-\tfrac32; &
\text{d) }&2^{2x+2}=2^{3}\Rightarrow 2x+2=3\Rightarrow x=\tfrac12.
\end{align*}
En c) on a utilisé $9^{x}=3^{2x}$ et $\tfrac{1}{27}=3^{-3}$ ; en d) $\tfrac{2^{2x+1}}{2^{x}}=2^{x+1}$ puis $8=2^{3}$.
\end{corrige}

\begin{corrige}[monotonie et inéquations]
\begin{itemize}[leftmargin=1.4em,topsep=1pt,itemsep=1pt]
  \item a) base $2>1$ (croissante) : $2^{x}<2^{5}\Leftrightarrow x<5$, soit $S=\,]-\infty,5[$.
  \item b) base $\tfrac12<1$ (décroissante) : le sens \emph{s'inverse}, $x<3$, soit $S=\,]-\infty,3[$.
  \item c) $2^{-1}<2^{0}<2^{1}<2^{3}$ (l'exposant croît, la fonction est croissante).
  \item d) base $\tfrac13<1$ : $\bigl(\tfrac13\bigr)^{2x}\leq\bigl(\tfrac13\bigr)^{6}\Leftrightarrow 2x\geq 6\Leftrightarrow x\geq 3$, soit $S=[3,+\infty[$.
\end{itemize}
\end{corrige}

\end{document}
